%%%%%%%%%%  scan idx 91  |  folio 78  %%%%%%%%%%
\SGAPrintPage{78}%
\sgaStableTarget{sga6.e01.expose}\section*{\underline{\underline{GENERALITIES ON FINITENESS CONDITIONS IN}}\\[3pt] \underline{\underline{DERIVED CATEGORIES}}}
\begin{center}
\underline{by L. ILLUSIE}
\end{center}

\sgaStableTarget{sga6.e01.sec.0}\subsection*{0. Introduction}
The purpose of Exposés I to IV is to develop, with the generality appropriate to this Seminar, the formalism of finiteness conditions in the derived categories of ringed topoi.

As was said in Exp.\ 0, it was the need to define, on arbitrary schemes, ``Grothendieck groups'' with good variance properties that required us to generalize and relax the notions of finiteness used until now.

A classical notion such as that of a coherent sheaf on a ringed space \((X,\underline{O}_X)\) becomes uninteresting as soon as \(\underline{O}_X\) ceases to be coherent. One might think of replacing it by the notion of finite presentation, but this has the drawback that the kernel of an epimorphism of Modules of finite presentation is in general no longer of finite presentation. We arrive at a satisfactory notion by observing that, when \(\underline{O}_X\) is coherent, a coherent sheaf \(F\) is not only of finite presentation but of \(n\)-finite presentation for every \(n\in\mathbb N\), ``\(n\)-finite presentation'' meaning that there locally exists an exact sequence

%%%%%%%%%%  scan idx 92  |  folio 79  %%%%%%%%%%
\SGAPrintPage{79}%
\[
  L_n\longrightarrow L_{n-1}\longrightarrow \ldots \longrightarrow L_0\longrightarrow F\longrightarrow 0,
\]
where the \(L_i\) are free of finite type. With the sheaf \(\underline{O}_X\) no longer assumed coherent, let us say that an \(\underline{O}_X\)-Module \(F\) is \uline{pseudo-coherent} when it satisfies the preceding condition, i.e.\ is of \(n\)-finite presentation for every \(n\). Then the notion of pseudo-coherence has, with respect to short exact sequences, the same stability property as coherence: when two of the terms of a short exact sequence of \(\underline{O}_X\)-Modules are pseudo-coherent, the third is too. We readily generalize the preceding notion to complexes of sheaves. A \uline{complex} \(F\) of \(\underline{O}_X\)-Modules is called \uline{pseudo-coherent} when it is \(n\)-pseudo-coherent for every \(n\in\mathbb Z\), ``\(n\)-pseudo-coherent'' meaning that there locally exists a quasi-isomorphism \(L\to F\), where \(L\) is a complex with degrees bounded above and components free of finite type in degree \(\geq n\). For a complex \(F\) concentrated in degree \(0\), it amounts to the same thing to say that \(F\) is pseudo-coherent as a complex or as a Module. The notion of a pseudo-coherent complex has excellent stability properties, described in Exp.\ I \S\S~1 and 2. First, it is stable under isomorphism in the derived category \(\Db(X)\) of the category of \(\underline{O}_X\)-Modules; moreover, when two of the objects of a distinguished triangle of \(\Db(X)\) are pseudo-coherent, the third is too: in other words, the full subcategory \(\Db(X)_{\mathrm{coh}}\) of \(\Db(X)\) formed by pseudo-coherent complexes is a triangulated subcategory. Furthermore, pseudo-coherence is preserved by derived inverse image and tensor product. When \(\underline{O}_X\) is coherent, saying that a complex \(F\) is pseudo-coherent simply means that \(F\) has locally bounded-above cohomology and that all the \(\underline H^i(F)\) are coherent.

When \(\underline{O}_X\) is a sheaf of regular local rings, every \(\underline{O}_X\)-Module
%%%%%%%%%%  scan idx 93  |  folio 80  %%%%%%%%%%
\SGAPrintPage{80}%
\(F\) that is coherent locally admits a \uline{finite} resolution to the left by free Modules of finite type, i.e.\ there locally exists an exact sequence
\[
  0\longrightarrow L_n\longrightarrow \ldots \longrightarrow L_0\longrightarrow F\longrightarrow 0,
\]
where the \(L_i\) are free of finite type. It is therefore natural, when \(\underline{O}_X\) is an arbitrary sheaf of local rings, to consider Modules with the preceding property. Such Modules are called \uline{perfect}. More generally, a \uline{complex} \(F\) of \(\underline{O}_X\)-Modules is called \uline{perfect} when there locally exists a quasi-isomorphism \(L\to F\), where \(L\) is a complex with bounded degrees and components free of finite type. It is shown in Exp.\ I \S~3 that the full subcategory \(\Db(X)_{\mathrm{parf}}\) of \(\Db(X)\) formed by perfect complexes is, like \(\Db(X)_{\mathrm{coh}}\), a triangulated subcategory, ``stable'' under derived inverse image and tensor product. We evidently have \(\Db(X)_{\mathrm{parf}}\subset\Db(X)_{\mathrm{coh}}\), but the inclusion is strict in general. Pseudo-coherence can be redefined ``by passage to the limit'' from perfection: a complex \(F\) is pseudo-coherent precisely when, for every \(n\in\mathbb Z\), \(F\) can be ``locally approximated to order \(n\) by a perfect complex'', by which we mean that there locally exists a distinguished triangle
\[
\begin{tikzcd}
& R \arrow[dl] &\\
L \arrow[rr] && F \arrow[ul] ,
\end{tikzcd}
\]
where \(L\) is perfect and \(R\) acyclic in degree \(\geq n\). Conversely, perfection is recovered from pseudo-coherence through an additional regularity condition: a complex is perfect precisely when it is pseudo-coherent and locally of finite tor-dimension (Exp.\ I \S~5). We thus recover the fact that, when the local rings are regular, every coherent sheaf is perfect, and, more generally, every pseudo-coherent complex with locally bounded cohomology is perfect.


%%%%%%%%%%  scan idx 94  |  folio 81  %%%%%%%%%%
\SGAPrintPage{81}%
Pseudo-coherence and perfection are the two fundamental notions of finiteness with which we shall work in this Seminar. Sections~1 to 5 of Exposé I are devoted to their definition and to establishing their elementary stability properties. For this, we use only two or three local properties of the category of free Modules of finite type with respect to the category of all Modules. It therefore proves convenient---and useful, particularly in Exp.\ II---to axiomatize the situation by introducing a notion of pseudo-coherence (resp.\ perfection) in a category fibered over a site relative to a fibered subcategory satisfying suitable conditions.

Sections~6 to 8 of Exp.\ I generalize to perfect complexes a number of familiar notions for locally free sheaves of finite type: rank, duality, trace of an endomorphism.

In Exp.\ II, we examine the problem of ``the existence of global resolutions'': under what conditions, on a ringed topos \((X,\underline{O}_X)\), is a pseudo-coherent (resp.\ perfect) complex \uline{globally} isomorphic in \(\Db(X)\) to a complex of locally free Modules of finite type? The importance of this question lies essentially in the fact that, at present, we do not know how to generalize to perfect complexes certain usual constructions on locally free Modules: exterior and symmetric powers, Chern classes (for more details on this subject, see Exp.\ XIV). This is why it is convenient to have workable sufficient criteria for the existence of global resolutions, which allow us to reduce certain questions about pseudo-coherent (resp.\ perfect) complexes to analogous questions about locally free sheaves of finite type. Such criteria
%%%%%%%%%%  scan idx 95  |  folio 82  %%%%%%%%%%
\SGAPrintPage{82}%
are given in Exp.\ II. In particular, it is shown that global resolutions exist in the following cases: \(X\) is the Zariski topos of an affine scheme, or more generally of a quasi-compact scheme with an ample invertible sheaf, or even an ample family of invertible sheaves in Kleiman's sense (for example a regular scheme); or again \(X\) is the topos of sheaves of sets on a compact topological space, \(\underline{O}_X\) the sheaf of continuous complex-valued functions. As an illustration---and to bring out the flexibility of the notions introduced---the appendix gives a ``purely sheaf-theoretic'' definition of the index of a family of elliptic operators.

Exposé III studies the stability of finiteness conditions under derived direct image. To obtain reasonable statements, we must ``relativize'' pseudo-coherence and perfection. We work here in the setting of ordinary schemes, which suffices for the Seminar, but there is no doubt that sooner or later an analogous theory will have to be developed for relative analytic schemes or spaces. Let \(p:X\to S\) be an \(S\)-scheme locally of finite type and \(F\) a complex of \(\underline{O}_X\)-Modules; we say that \(F\) is \uline{pseudo-coherent} (\uline{resp.\ perfect}) \uline{relative to} \(p\) when we can locally embed (by a closed immersion) \(X\) in a smooth \(S\)-scheme \(X'\) so that \(F\), extended by zero to \(X'\), is pseudo-coherent (resp.\ perfect). Pseudo-coherence relative to \(p\) is especially interesting when \(S\) is not locally noetherian, since in the opposite case it coincides with ordinary pseudo-coherence. Moreover, as one might expect, saying that \(F\) is perfect rel.\ to \(p\)
%%%%%%%%%%  scan idx 96  |  folio 83  %%%%%%%%%%
\SGAPrintPage{83}%
amounts to saying that \(F\) is pseudo-coherent rel.\ to \(p\) and locally of finite tor-dimension rel.\ to \(p\) (i.e.\ rel.\ to the sheaf of rings \(p^{-1}(\underline{O}_S)\)). The central theorem of Exp.\ III is the \uline{finiteness theorem}, which asserts (in a slightly more precise form) that, for a \uline{proper} morphism \(f:X\to Y\) of \(S\)-schemes locally of finite type, the functor \(\mathbb{R}f_*\) sends complexes pseudo-coherent rel.\ to \(S\) to complexes pseudo-coherent rel.\ to \(S\); this theorem is in fact a conjecture, but is nevertheless proved in the following two special cases: a)~\(S\) is locally noetherian b)~\(f\) is projective; unfortunately, it seems that extending it to the general case is of the same order of difficulty as the analogous theorem in Analytic Geometry (Grauert's ``theorem''). Combining the finiteness theorem with an essentially trivial formula called the \uline{projection formula}, we obtain workable criteria for the stability of relative perfection under direct image. As a corollary, we recover ``Grauert's continuity and semicontinuity theorems'' (EGA III 7.6).

Exposé IV translates the results of Exposés I to III into the language of ``Grothendieck groups''. On a ringed topos \((X,\underline{O}_X)\), we denote by \(K^\bullet(X)\) (resp.\ \(K_\bullet(X)\)) the ``Grothendieck group'' of the category of complexes that are perfect and of finite tor-dimension (resp.\ pseudo-coherent and with bounded cohomology). The group \(K^\bullet(X)\) is a ring, and \(K_\bullet(X)\) a module over \(K^\bullet(X)\). \(K^\bullet\) is a contravariant functor (as the raised dot suggests), while \(K_\bullet\) is a covariant functor for proper pseudo-coherent morphisms (of schemes or analytic spaces, and subject to having the finiteness theorem\dots). There are also unusual variances,
%%%%%%%%%%  scan idx 97  |  folio 84  %%%%%%%%%%
\SGAPrintPage{84}%
covariance of \(K^\bullet\), contravariance of \(K_\bullet\), for morphisms satisfying suitable regularity hypotheses. Finally, the globalization criteria of Exp.\ II allow us to relate these to the ``Grothendieck groups'' defined naively from coherent sheaves or locally free sheaves of finite type.

\sgaStableTarget{sga6.e01.sec.1}\section*{1. Preliminary definitions}

\sgaStableTarget{sga6.e01.sec.1.1}\subsection*{1.1. Fibered categories with additive (resp. abelian, resp. triangulated) fibres}

\sgaStableTarget{sga6.e01.para.1-1-1}1.1.1. Let \(\underline{S}\) be a category, and \(\underline{C}\) a fibered \(\underline{S}\)-category (SGA 1 VI 6.1) whose fibres are additive (resp. triangulated) categories. We shall say that \(\underline{C}\) is an \uline{additive} \(\underline{S}\)\uline{-category} (\uline{resp.\ triangulated}) (or that \(\underline{C}\) is additive (resp.\ triangulated) over \(\underline{S}\)) when, for every arrow \(f:X\to Y\) of \(\underline{S}\), the inverse-image functor (determined up to unique isomorphism) \(f^*:\underline{C}_Y\longrightarrow \underline{C}_X\) is additive (resp.\ exact). We shall say that \(\underline{C}\) is an \uline{abelian} \(\underline{S}\)-category when \(\underline{C}\) is an additive \(\underline{S}\)-category with abelian fibres, and that \(\underline{C}\) is a \uline{flat} abelian \(\underline{S}\)-category when, moreover, the inverse-image functors are exact.

\sgaStableTarget{sga6.e01.para.1-1-2}1.1.2. Let \(\underline{C}\) be a triangulated \(\underline{S}\)-category. We define, for example using a cleavage of \(\underline{C}\) (SGA 1 VI 7.1), a fibered \(\underline{S}\)-category \(\Trc(\underline{C})\) whose fibre at each \(X\in \ob(\underline{S})\) is the category \(\Trc(\underline{C}_X)\) of distinguished triangles of \(\underline{C}_X\).

\sgaStableTarget{sga6.e01.para.1-1-3}1.1.3. Let \(\underline{C}\) and \(\underline{C}'\) be additive (resp. triangulated) \(\underline{S}\)-categories and \(F:\underline{C}\to\underline{C}'\) a functor. We call \(F\) an \uline{additive} (resp.\ \uline{exact}) \(\underline{S}\)\uline{-functor} when \(F\) is a cartesian \(\underline{S}\)-functor that is additive (resp.\ exact) on each fibre.
% Additional macro definitions
\providecommand{\cC}{\mathcal C}
\providecommand{\cD}{\mathcal D}
\providecommand{\cK}{\mathcal K}
\providecommand{\cO}{\mathcal O}
\providecommand{\cS}{\mathcal S}
\providecommand{\Cpx}{\mathrm C}
\providecommand{\Kpx}{\mathrm K}
\providecommand{\Dpx}{\mathrm D}
\providecommand{\D}{\mathrm D}
\providecommand{\Ho}{\mathrm H}
\providecommand{\RHom}{\mathrm R\!\operatorname{Hom}}
\providecommand{\Ob}{\operatorname{ob}}
\providecommand{\Coker}{\operatorname{Coker}}
\providecommand{\Ker}{\operatorname{Ker}}
\renewcommand{\Im}{\operatorname{Im}}
\providecommand{\Hom}{\operatorname{Hom}}
\providecommand{\Cone}{\operatorname{Cone}}
\providecommand{\id}{\operatorname{Id}}
\providecommand{\qis}{\xrightarrow{\sim}}
%%%%%%%%%%  scan idx 98  |  folio 85  %%%%%%%%%%
\SGAPrintPage{85}%
\sgaStableTarget{sga6.e01.para.1-1-3.occ2}When \(\underline{C}\) and \(\underline{C}'\) are triangulated \(\underline{S}\)-categories and \(F\) is an exact \(\underline{S}\)-functor, \(F\) defines a cartesian \(\underline{S}\)-functor \(\Trc(\underline{C})\longrightarrow \Trc(\underline{C}')\).

\sgaStableTarget{sga6.e01.para.1-1-4}1.1.4. Let \(\underline{C}\) be a triangulated \(\underline{S}\)-category, and \(\underline{C}'\) a full fibered sub-\(\underline{S}\)-category of \(\underline{C}\). We call \(\underline{C}'\) a \uline{triangulated sub-}\(\underline{S}\)\uline{-category} of \(\underline{C}\) when, for every \(X\in\Ob \underline{S}\), \(\underline{C}'_X\) is a triangulated subcategory of \(\underline{C}_X\) ([V] chap.\ I \S~1 2.3).

\sgaStableTarget{sga6.e01.para.1-1-5}1.1.5. Let \(\underline{C}\) be an additive \(\underline{S}\)-category. We define (SGA 4 XVII 2.2, 2.4) an additive (resp.\ triangulated) \(\underline{S}\)-category \(\Cpx(\underline{C})\), (resp.\ \(\Kpx(\underline{C})\)), whose fibre at \(X\in\Ob \underline{S}\) is the category \(\Cpx(\underline{C}_X)\) (resp.\ \(\Kpx(\underline{C}_X)\)) of complexes\SGAAuthorFootnote{1}{with differential of degree \(+1\).} (resp.\ complexes ``up to homotopy'') of \(\underline{C}_X\). When \(\underline{C}\) is a flat abelian \(\underline{S}\)-category, we define (loc.\ cit.) a triangulated \(\underline{S}\)-category \(\Dpx(\underline{C})\), whose fibre at \(X\in\Ob \underline{S}\) is the derived category \(\Dpx(\underline{C}_X)\); for an arrow \(f:X\to Y\) of \(\underline{S}\), we shall denote by \(f^*:\Dpx(\underline{C}_Y)\longrightarrow \Dpx(\underline{C}_X)\) the inverse-image functor, obtained from \(f^*:\Kpx(\underline{C}_Y)\longrightarrow \Kpx(\underline{C}_X)\) by simple passage to the quotient\SGAAuthorFootnote{2}{for \(a\in\Ob \Dpx(\underline{C}_Y)\) (resp.\ \(\mathrm{Fl}\,\Dpx(\underline{C}_Y)\)), we shall often write, in abbreviated form, \(a_X\) instead of \(f^*(a)\).}. We define full triangulated sub-\(\underline{S}\)-categories \(\Kpx^*(\underline{C})\) (resp.\ \(\Dpx^*(\underline{C})\)) \((*=+,-,b)\) by the usual degree conditions.

\sgaStableTarget{sga6.e01.sec.1.2}1.2. \uline{Terminology.} Let \(\underline{S}\) be a site, \(\underline{C}\) a flat abelian \(\underline{S}\)-category (1.1), and \(\underline{C}_0\) a strictly full fibered sub-\(\underline{S}\)-category.

Let \(F\in\Ob\underline{C}_X\), for \(X\in\Ob \underline{S}\). The set of \(f:U\to X\) in \(\underline{S}_{/X}\) such that there exists an epimorphism \(E\longrightarrow f^*F\) with \(E\in\Ob\underline{C}_{0U}\) is a sieve on \(X\). When it is covering, we shall say that \(F\) is of \(\underline{C}_0\)-\uline{finite type} (or simply of finite type when there is no risk of confusion).


